\section{独立数值与采样分布验证}

最终测试日志中的 335 项测试通过，另有 28 个参数子场景。新增核验同时比较能量恒等式和采样分布；后者直接检验参考势变化后，残差校正是否恢复目标分布。

\begin{center}\small
\begin{tabular}{@{}lr@{}}\toprule
验证对象 & 最大绝对误差\\\midrule
半谱总能 / 独立全谱 & $5.912\times10^{-12}$ K\\
增量能量 / 独立有限和 & $8.640\times10^{-12}$ K\\
框架 FFT 节点 / 直接求和 & $1.273\times10^{-11}$ K\\
平移相位 / 坐标重算结构因子 & $7.845\times10^{-15}$\\
Boltzmann 插值 / 独立加权配分和 & $1.990\times10^{-13}$ K\\
生产目标探针 / 独立 Ewald & $9.095\times10^{-12}$ K\\
\bottomrule\end{tabular}\end{center}

非均匀根位点检验把 85\% 的提议概率集中在 24 个 voxel 中的一个，同时保留全空间支持。理论计数均值和方差均为 2.7；16,000 个保存样本给出均值 2.71625、方差 2.73536。该算例检验根位点密度只进入 Rosenbluth 权重，而没有重复进入残差。

更强的检验故意把参考势截幅到 50 K，使其与带电单粒子的完整目标明显不同。独立空间积分和原生采样比较如下；误差条为分块标准误。

\begin{center}\small
\begin{tabular}{@{}lrrr@{}}\toprule
观测量 & 独立目标积分 & 未校正参考势 & 校正后采样\\\midrule
$\langle \cos(2\pi x/L_x)\rangle$ & -0.024514 & -0.001467 & $-0.013031\pm0.010045$\\
$\langle \sin(2\pi x/L_x)\rangle$ & -0.113245 & -0.007541 & $-0.109579\pm0.009006$\\
$\langle \cos(2\pi z/L_z)\rangle$ & 0.011903 & 0.003535 & $0.009988\pm0.009983$\\
\bottomrule\end{tabular}\end{center}

独立积分的网格加密最大变化为 $6.119\times10^{-5}$；本次实际残差范围为 -570.50 至 382.85 K。参考势本身不能复现目标分布，加入第二阶段后采样与目标积分一致。逐分子结构因子恢复、五类移动和缓存跨温度重用也已核验。

\section{完整循环的同机单线程基准}

\subsection{模型与计时约定}

CPU 为 AMD EPYC 9V74 80-Core Processor。两端使用相同编译优化，固定一个 OpenMP / BLAS 线程。对照源为打包的未修改 0.5 核心；0.6 所有主表均保持同一完整矩形有限 Ewald 目标，没有施加径向模式截断。

本轮没有用户实际 CIF 与完整 p-xylene 力场输入，因此创建并公开了计算规模匹配的合成算例：48 \AA{} 立方盒、4,416 个框架点、12 个客体位点、10 个 CBMC trial、10 个初始分子、$k_{\max}=13$、$T=350$ K、$r_c=10$ \AA{}、$\alpha=0.28$ \AA$^{-1}$。五种移动等概率，平移上限 0.65 \AA{}、转角上限 0.35 rad；LJ 混合采用 Lorentz--Berthelot 与截断移位，1--4 LJ / Coulomb 缩放为 0.5 / 0.75。

框架为形成周期狭缝的八个轻微起伏平面，位点 $\sigma=2.9$ \AA、$\epsilon=4$ K；电荷由固定种子生成并减去均值。刚性算例由六个环 C、四个 H 和两个甲基伪位点组成，采用公开在脚本中的合成参数；柔性算例为十二位点链，$\sigma=2.2$ \AA、$\epsilon=6$ K，电荷交替 $\pm0.03e$，键长 1.3 \AA、键角 $112^\circ$、三重扭转系数 15 K。两类活度分别为 $8/V$ 和 $12/V$。这些参数定义数值和性能任务，不是对某种材料的力场标定。

刚性每组预热 500 次、生产 3000 次；柔性分别为 200 与 1500 次。各三种子、独立进程并交替测试顺序。生产吞吐不含初始化和末态参考重算，关闭逐帧总能记录；接受次数与计数轨迹同时保留。

\subsection{十二位点刚性客体}

\begin{center}\small
\begin{tabular}{@{}lrrrr@{}}\toprule
路径 & 尝试/s & 接受/s & 配对加速 & 加速范围\\\midrule
0.5 全谱逐候选 & 316.10 & 197.14 & 1.00 & 1.00--1.00\\
0.6 半谱逐候选 & 523.38 & 326.42 & 1.67 & 1.66--1.70\\
0.6 参考网格 $32^3$ & 4,933.15 & 2,428.76 & 16.71 & 15.53--17.14\\
$32^3$ + alias 根位点 & 4,972.70 & 2,405.82 & 16.41 & 15.73--16.79\\
0.6 参考网格 $64^3$ & 4,499.63 & 2,507.79 & 14.23 & 11.27--16.23\\
\bottomrule\end{tabular}\end{center}

\subsection{十二位点柔性客体}

\begin{center}\small
\begin{tabular}{@{}lrrrr@{}}\toprule
路径 & 尝试/s & 接受/s & 配对加速 & 加速范围\\\midrule
0.5 全谱逐候选 & 101.32 & 63.83 & 1.00 & 1.00--1.00\\
0.6 半谱逐候选 & 186.20 & 116.90 & 1.85 & 1.83--1.92\\
0.6 参考网格 $32^3$ & 4,683.53 & 2,740.62 & 46.22 & 43.00--49.43\\
\bottomrule\end{tabular}\end{center}

吞吐两列分别为三次测量的中位数；加速列为每次配对比值的中位数，因此不必等于两列中位数之比。接受吞吐避免把纯拒绝路径的高尝试数误认成构型探索效率；构型相关性另在长链部分检验。

全部配对任务的末态完整重算最大差为 $9.286\times10^{-9}$ K。不同接受算法使用不同数量的随机数，轨迹无需逐步相同；物理输入哈希、初始能量、末态重算和分布共同定义比较依据。

\subsection{缓存初始化与生产工作量}

\begin{center}\small
\begin{tabular}{@{}lrrr@{}}\toprule
网格 & 首次构建/s & 读取复用/s & 磁盘字段/MiB\\\midrule
$32^3$ & 0.672 & 0.115 & 1.13\\
$64^3$ & 4.689 & 0.282 & 9.00\\
\bottomrule\end{tabular}\end{center}

这里的初始化包括势函数和原生引擎构建，单独记录，不加入生产吞吐。缓存保存原始字段，重新选择温度、活度和根位点偏置时可以继续读取。

\begin{center}\small
\begin{tabular}{@{}lrrr@{}}\toprule
路径 & 第一阶段通过/\% & 第二阶段通过/\% & 精确分子评价/尝试\\\midrule
0.6 参考网格 $32^3$ & 67.95 & 72.50 & 1.109\\
$32^3$ + alias 根位点 & 67.93 & 70.50 & 1.105\\
0.6 参考网格 $64^3$ & 66.69 & 89.07 & 1.089\\
\bottomrule\end{tabular}\end{center}

阶段统计包含预热与生产；空体系无法执行的删除等移动不进入第一阶段分母。一次替换包含旧、新两个残差，旧分子结构因子与两种局域能直接读取缓存。

\section{粗网格的负结果与混合检验}

初始三线性参考把非常高的排斥节点能量混入相邻可达体积。下表给出相同刚性任务三条短链的平均计数，保留失败的优化路径以解释最终设计。Boltzmann 插值恢复了进入近壁区域的机会，但参数选择仍以更长轨迹和相关性判断。

\begin{center}\small
\begin{tabular}{@{}lr@{}}\toprule
路径 & 三种子的平均粒子数\\\midrule
精确 0.5 参考 & 13.79, 11.97, 14.24\\
线性能量插值 $32^3$ & 9.45, 8.86, 8.63\\
Boltzmann $32^3$ & 11.44, 13.36, 10.58\\
Boltzmann $64^3$ & 11.10, 13.70, 13.30\\
\bottomrule\end{tabular}\end{center}

长链对比每种路径运行三个独立种子，每条预热 5,000 次、生产 30,000 次，每 10 次保存计数。下表的误差为三个种子均值的标准误；ESS 按各链的初始正序列自相关估计后相加，仅针对粒子数这一观测量。

\begin{center}\small
\begin{tabular}{@{}lrrr@{}}\toprule
路径 & $\langle N\rangle$ & 种子均值 SE & $N$ 的 ESS 合计\\\midrule
0.5 全谱逐候选 & 13.936 & 0.395 & 199.6\\
0.6 参考网格 $32^3$ & 13.799 & 0.708 & 119.4\\
0.6 参考网格 $64^3$ & 13.242 & 0.243 & 220.8\\
$32^3$ + alias 根位点 & 14.451 & 0.356 & 129.8\\
\bottomrule\end{tabular}\end{center}

相同 90,000 次生产尝试下，$64^3$ 路径的计数 ESS 与精确对照相近，且优于本例的 $32^3$ 路径；结合完整循环吞吐，本报告建议从 $64^3$ Boltzmann 参考网格开始。$32^3$ 与 alias 参数仍可按目标体系实测选择。独立链均值仍有散布，表中没有把这组诊断作为精确等温线的最终物性。长链为缩短验证时间并行运行，其耗时不用于性能倍率或 ESS/s。

\section{显式模式截断的精度与成本空间}

除保持原有限和的半谱优化外，本版提供离线误差界建议工具。给定最大环境载量与完整分子的电荷绝对值和，可以界定省略倒空间壳层造成的能量增量差。下面使用同一刚性合成模型、$N_{\max}=64$；保守替换界覆盖完整分子的插入、删除与替换。主性能表没有采用这些额外截断。

\begin{center}\small
\begin{tabular}{@{}lrrrr@{}}\toprule
容限/K & $k_c$/\AA$^{-1}$ & 保留半谱 & 替换上界/K & 实测最大差/K\\\midrule
0.01 & 2.306566 & 8851 & 0.009744 & $1.162\times10^{-5}$\\
0.1 & 2.168751 & 8193 & 0.097322 & $7.232\times10^{-5}$\\
1 & 2.038423 & 7289 & 0.973553 & $6.917\times10^{-4}$\\
\bottomrule\end{tabular}\end{center}

每个建议另以环境占据 0、10、64 的 12 组新旧刚体构型检查实际省略差。解析界是给定载量范围内的约束，样本最大差只是这组构型的测量。非零径向截断定义另一个有限 Ewald Hamiltonian；没有把其模式数缩减当作本轮已测得的额外加速。

\begin{lstlisting}
python tools06/advise_reciprocal.py your_hpc_input.json \
    --max-molecules 64 --tolerance-K 0.1 \
    --output reciprocal_advice.json
\end{lstlisting}

该工具只输出建议，不修改输入文件。载量超出分析范围时应重新评估，生产采样中不动态切换目标模式集合。
